\setlength{\absparsep}{18pt} % Espaçamento entre parágrafos
\begin{resumo}
 
A produção de biodiesel por transesterificação alcalina enfrenta dificuldades relacionadas à baixa miscibilidade entre o triglicerídeo e o álcool utilizado na reação, o que tende a comprometer a eficiência do processo e a etapa de separação de fases subsequente. Nesse contexto, líquidos iônicos próticos (LIPs) têm sido investigados como cossolventes capazes de modular o comportamento de equilíbrio líquido-líquido (ELL) desses sistemas, ainda que os mecanismos pelos quais essa modulação ocorre permaneçam pouco explorados na literatura, sobretudo em sistemas ternários envolvendo triglicerídeos, para os quais dados experimentais de ELL não estão disponíveis. Este trabalho teve como objetivo modelar termodinamicamente o ELL de doze sistemas ternários a 298,15~K, combinando três triglicerídeos com grau crescente de insaturação (trioleína, trilinoleína e trilinolenina), dois álcoois (etanol e metanol) e dois LIPs derivados do ânion hexanoato, o hexanoato de 2-hidroxietilamônio ([2-HEA][Hx]) e o de dietilenotriamônio ([DETA][Hx]), avaliando por duas abordagens preditivas de natureza distinta se esses compostos aumentam ou reduzem a miscibilidade do par óleo-álcool: o UNIFAC-LL, modelo de contribuição de grupos, e o COSMO-SAC, modelo \textit{a priori} baseado em perfis de densidade de carga superficial obtidos por cálculo quântico. Os parâmetros de interação ausentes na matriz do UNIFAC-LL para os grupos característicos dos LIPs foram preenchidos por um modelo de aprendizado de máquina, e o motor de cálculo comum aos dois modelos foi verificado externamente contra sistemas de referência da literatura. Nos dezoito sistemas para os quais se obteve curva binodal, sem exceção, a largura da lacuna de imiscibilidade aumenta com a fração molar de LIP, com inclinações entre $+0{,}87$ e $+2{,}95$, de modo que a hipótese de cossolvência não se confirma nas condições modeladas: o LIP afasta as fases coexistentes em vez de aproximá-las, comportando-se como agente de separação. O alargamento decorre, em ambos os modelos, do afastamento do ramo rico em óleo, ao passo que a composição da fase alcoólica permanece praticamente inalterada, o que sugere que os ganhos de rendimento reportados experimentalmente, obtidos em temperatura e concentração distintas, não se explicam por solubilização em fase homogênea. Os modelos concordam quanto aos sistemas com [2-HEA][Hx], para os quais ambos indicam região bifásica, e divergem quanto aos sistemas com [DETA][Hx], em que o COSMO-SAC prevê separação de fases nos seis casos e o UNIFAC-LL não permite construir a curva em nenhum deles, divergência que os dados obtidos não permitem arbitrar na ausência de medidas experimentais. A interpretação físico-química apoia-se no ancoramento do LIP à fase alcoólica por ligação de hidrogênio, na competição resultante pela hidroxila do álcool e no impedimento estérico das cadeias esterificadas ao glicerol.
 
\vspace{\onelineskip}
\noindent
\textbf{Palavras-chave}: Líquidos iônicos próticos. Biodiesel. Equilíbrio líquido-líquido. UNIFAC-LL. COSMO-SAC.
 
\end{resumo}